\section*{Capítulo \ref{ch:b2:metamorphic-facies} --- Reações e fácies metamórficas}

\begin{solution}{exo:b2:metamorphic-facies:1}
(a) \omterm{def:b2:metamorphic-facies:metamorphism}{Metamorfismo} de contato, um \omterm{def:b2:metamorphic-facies:metamorphism}{metamorfismo} térmico local pelo calor da intrusão.
(b) \omterm{def:b2:metamorphic-facies:metamorphism}{Metamorfismo} de deslocamento, em uma zona de falha: redução do tamanho dos grãos, um milonito.
(c) \omterm{def:b2:metamorphic-facies:metamorphism}{Metamorfismo} de fundo oceânico, pela água do mar quente perto da dorsal, com
recristalização incompleta. (d) \omterm{def:b2:metamorphic-facies:metamorphism}{Metamorfismo} de soterramento, regional, sem deformação nem
xistosidade.
\end{solution}

\begin{solution}{exo:b2:metamorphic-facies:2}
$\Delta_r S = 95.40 - 82.80 = \qty{12.60}{J.K^{-1}}$, $\Delta_r V = 4.986 - 4.415 =
\qty{0.571}{J.bar^{-1}}$; inclinação $12.60/0.571 = \qty{22.1}{bar/K} = \qty{2.21}{MPa/K}$.
Em \qty{100}{\celsius}, a pressão de equilíbrio sobe \qty{2.21}{kbar}.
\end{solution}

\begin{solution}{exo:b2:metamorphic-facies:3}
(a) Xisto verde. (b) Xisto azul. (c) Anfibolito: hornblenda com plagioclásio de mais
de \qty{17}{\percent} de anortita. (d) Eclogito.
\end{solution}

\begin{solution}{exo:b2:metamorphic-facies:4}
(a) $\rho = M/V$: cianita $162.046/44.15 = \qty{3.670}{g.cm^{-3}}$, andaluzita
$162.046/51.52 = 3.145$, silimanita $162.046/49.86 = 3.250$. (b) A cianita: a pressão mais
alta, o termo $V\,(P - 1)$ da energia de Gibbs cresce menos para o menor volume.
(c) Os cristais naturais contêm outros elementos (o ferro férrico, por exemplo, substitui parte do
alumínio na andaluzita e na silimanita naturais), e as medidas em
espécimes naturais são menos precisas que os volumes molares das tabelas.
\end{solution}

\begin{solution}{exo:b2:metamorphic-facies:5}
(a) $\Delta_r G = 3900 - 298.15\times8.59 = \qty{+1339}{J} > 0$: cianita. (b) A
\qty{1}{bar}, $\Delta_r H = T\Delta_r S$: $T = 3900/8.59 = \qty{454}{K} =
\qty{181}{\celsius}$. (c) $P = (8.59\times673.15 - 3900)/0.737 = 1882/0.737 =
\qty{2554}{bar} = \qty{2.55}{kbar}$, contra 2.54 medidos. (d) A reação é lenta demais
a baixa temperatura para avançar: a andaluzita sobrevive de forma metaestável, como a maioria dos minerais
de alta temperatura na superfície.
\end{solution}

\begin{solution}{exo:b2:metamorphic-facies:6}
(a) $\Delta_r H = 7700 + 1000\times(121.39 - 122.29) = \qty{6800}{J}$; $\Delta_r S = 296.44
- 285.02 = \qty{11.42}{J.K^{-1}}$; $\Delta_r G = 6800 - 1000\times11.42 = \qty{-4620}{J}$.
(b) Com os valores a \qty{298.15}{K}: $7700 - 1000\times12.60 = \qty{-4900}{J}$.
(c) A silimanita. O método simples erra em \qty{280}{J}, cerca de \qty{6}{\percent}; dividido
por $\Delta_r V = 0.571$, isso dá cerca de \qty{490}{bar} na linha, muito menos que o
efeito da incerteza das entalpias.
\end{solution}

\begin{solution}{exo:b2:metamorphic-facies:7}
(a) $\Delta_r G = -1127.4 + 1128.5 = \qty{1.1}{kJ}$ e $\Delta_r V = (34.15 -
36.93)/10 = \qty{-0.278}{J.bar^{-1}}$: $P = 1100/0.278 = \qty{3957}{bar} \approx
\qty{4.0}{kbar}$. (b) $\Delta_r S = 88.0 - 91.7 = -3.7$; inclinação $-3.7/-0.278 =
\qty{13.3}{bar/K}$. (c) Com $\Delta_r H = 0$, $P = T\Delta_r S/\Delta_r V = 13.31\,T$:
a \qty{573.15}{K}, \qty{7629}{bar} $= \qty{7.63}{kbar}$. (d) A aragonita das conchas é
metaestável na superfície: os organismos a produzem, e ela persiste porque só se transforma em
calcita lentamente.
\end{solution}

\begin{solution}{exo:b2:metamorphic-facies:8}
(a) $-32.4/-1.698 = \qty{19.1}{bar/K}$. (b) $P = (T\Delta_r S - \Delta_r H)/\Delta_r V =
(-32.4\times773.15 + 5000)/(-1.698) = (-25050 + 5000)/(-1.698) = \qty{11808}{bar}
\approx \qty{11.8}{kbar}$. (c) $\sqrt{2.6^2 + 3.6^2 + 1.0^2} = \qty{4.55}{kJ}$ sobre
$\Delta_r H$, logo $4550/1.698 = \qty{2.7}{kbar}$ sobre a pressão. (d) Às pressões da
fácies eclogito, o componente albita do plagioclásio é instável em relação à
jadeíta e ao quartzo; o sódio vai para o piroxênio (onfacita), e o plagioclásio
desaparece.
\end{solution}

\begin{solution}{exo:b2:metamorphic-facies:9}
(a) De 1 a \qty{2}{kbar}: $1000/34 = \qty{29.4}{bar/K}$; de 2 a 3: $1000/25 =
\qty{40.0}{bar/K}$. (b) Com $\Delta_r S$ fixo, $\Delta_r V$ é inversamente proporcional à
inclinação: diminuiu por um fator $40.0/29.4 = 1.36$. (c) A água liberada: seu volume molar
cai à medida que a pressão sobe, enquanto os volumes dos sólidos quase não mudam.
\end{solution}

\begin{solution}{exo:b2:metamorphic-facies:10}
(a) A \qty{873.15}{K}: cianita $=$ andaluzita $(8.59\times873.15 - 3900)/0.737 =
\qty{4885}{bar}$; cianita $=$ silimanita $(12.60\times873.15 - 7700)/0.571 =
\qty{5782}{bar}$; andaluzita $=$ silimanita $(4.01\times873.15 - 3800)/(-0.166) =
\qty{1799}{bar}$. (b) A \qty{4885}{bar}, abaixo da linha cianita $=$ silimanita (5782),
a silimanita é mais estável que a cianita, e acima da linha andaluzita $=$ silimanita
(1799), mais estável que a andaluzita: o ponto fica no campo da silimanita, e o
equilíbrio cianita $=$ andaluzita é ali metaestável. (c) A \qty{4}{kbar}: entre 1799
e \qty{5782}{bar}, silimanita. A \qty{7}{kbar}: acima de \qty{5782}{bar} a cianita vence a
silimanita, e acima de \qty{4885}{bar} vence a andaluzita: cianita.
\end{solution}

\begin{solution}{exo:b2:metamorphic-facies:11}
(a) $z = P/(\rho g) = 3.75\times10^8/(2700\times9.81) = \qty{14.2}{km}$, com incerteza de
$0.25\times10^8/26487 = \qty{0.94}{km}$. (b) $(504 - 10)/14.16 =
\qty{35}{\celsius.km^{-1}}$. (c) Os três minerais só coexistem em equilíbrio no
\omterm{def:b2:metamorphic-facies:invariant}{ponto invariante}: a trajetória da rocha precisa passar por ele, e as reações precisam
parar ali. A maioria das rochas passa ao lado dele e conserva um ou dois polimorfos.
\end{solution}

\begin{solution}{exo:b2:metamorphic-facies:12}
(a) X: $T = 10 + 10\times30 = \qty{310}{\celsius}$ (\qty{583.15}{K}); $P = 3000\times
9.81\times30000 = \qty{8.83e8}{Pa} = \qty{8.83}{kbar}$. Y: $T = 10 + 30\times20 =
\qty{610}{\celsius}$ (\qty{883.15}{K}); $P = 2800\times9.81\times20000 = \qty{5.49}{kbar}$.
(b) X: calcita $=$ aragonita a $13.31\times583.15 = \qty{7.76}{kbar} < 8.83$, aragonita;
albita $=$ jadeíta + quartzo a $(32.4\times583.15 - 5000)/1.698 = \qty{8.18}{kbar} <
8.83$, jadeíta + quartzo; cianita $=$ andaluzita a $(8.59\times583.15 - 3900)/0.737 =
\qty{1.51}{kbar}$, abaixo da temperatura do ponto triplo: cianita. Y: calcita $=$
aragonita a \qty{11.75}{kbar} $> 5.49$, calcita; albita $=$ jadeíta + quartzo a
\qty{13.9}{kbar}, albita; acima do ponto triplo, cianita $=$ silimanita a
\qty{6.00}{kbar} $> 5.49$ e andaluzita $=$ silimanita a \qty{1.56}{kbar} $< 5.49$:
silimanita. (c) X, a \qty{10}{\celsius\per\km}, é a trajetória mais fria: jadeíta e
aragonita a apenas \qty{310}{\celsius} a situam na fácies xisto azul. Y, três vezes
mais quente, termina com silimanita a \qty{610}{\celsius} e pressão moderada, na
fácies anfibolito.
\end{solution}

\begin{solution}{pb:b2:metamorphic-facies:1}
\noindent\textbf{Parte I.}

\textbf{1.} $\Delta_r H = -2589.9 + 2593.8 = \qty{3.9}{kJ}$; $\Delta_r S = 91.39 - 82.80 =
\qty{8.59}{J.K^{-1}}$; $\Delta_r V = 5.152 - 4.415 = \qty{0.737}{J.bar^{-1}}$.

\textbf{2.} $\qty{7.7}{kJ}$; $\qty{12.60}{J.K^{-1}}$; $\qty{0.571}{J.bar^{-1}}$.

\textbf{3.} $\qty{3.8}{kJ}$; $\qty{4.01}{J.K^{-1}}$; $\qty{-0.166}{J.bar^{-1}}$; cada um é
o segundo menos o primeiro, já que a andaluzita se cancela.

\textbf{4.} A silimanita tem a maior entropia (95.40): ela fica do lado de alta temperatura
de suas duas linhas. A cianita tem o menor volume (4.415): ela fica do lado de
alta pressão de suas duas linhas. A andaluzita, de maior volume, ocupa o
lado de baixa pressão.

\textbf{5.} $8.59/0.737 = \qty{11.66}{bar/K}$; $12.60/0.571 = \qty{22.07}{bar/K}$;
$4.01/(-0.166) = \qty{-24.16}{bar/K}$.

\textbf{6.} A silimanita tem mais entropia, mas menos volume que a andaluzita: $\Delta_r S > 0$
e $\Delta_r V < 0$. Elevar a temperatura favorece a silimanita, elevar a pressão
também, de modo que a linha vai para pressões mais baixas a temperaturas mais altas.

\noindent\textbf{Parte II.}

\textbf{7.} $3900 - 298.15\times8.59 = \qty{+1339}{J}$; $7700 - 298.15\times12.60 =
\qty{+3943}{J}$. Ambos positivos: a cianita é estável na superfície.

\textbf{8.} $T = \Delta_r H/\Delta_r S$: $3900/8.59 = \qty{454.0}{K} = \qty{181}{\celsius}$;
$7700/12.60 = \qty{611.1}{K} = \qty{338}{\celsius}$; $3800/4.01 = \qty{947.6}{K} =
\qty{674}{\celsius}$.

\textbf{9.} Aquecendo a \qty{1}{bar}, a cianita dá lugar à andaluzita a
\qty{181}{\celsius}; a \qty{338}{\celsius} a cianita já não é estável, de modo que cianita $=$
silimanita é ali metaestável; a andaluzita dá lugar à silimanita a
\qty{674}{\celsius}. O primeiro e o terceiro são equilíbrios estáveis.

\textbf{10.} $(8.59\times673.15 - 3900)/0.737 = \qty{2554}{bar} = \qty{2.55}{kbar}$,
a \qty{0.01}{kbar} do experimento.

\textbf{11.} $(12.60\times973.15 - 7700)/0.571 = 4562/0.571 = \qty{7989}{bar} =
\qty{7.99}{kbar}$, \qty{0.13}{kbar} acima do experimento.

\textbf{12.} Cianita $=$ andaluzita sobe a partir de \qty{181}{\celsius} a \qty{1}{bar} com
inclinação de \qty{11.7}{bar/K}; andaluzita $=$ silimanita desce de cerca de \qty{11.5}{kbar} a
\qty{200}{\celsius} até \qty{1}{bar} a \qty{674}{\celsius}; cianita $=$ silimanita, a
mais íngreme das duas linhas ascendentes, cruza as duas perto de \qty{514}{\celsius} e \qty{3.9}{kbar}, cada linha estável de
um só lado do cruzamento (o diagrama do capítulo).

\noindent\textbf{Parte III.}

\textbf{13.} $a_1 = 8.59/0.737 = \qty{11.655}{bar/K}$, $b_1 = 3900/0.737 =
\qty{5291.7}{bar}$; $a_2 = 12.60/0.571 = \qty{22.067}{bar/K}$, $b_2 = 7700/0.571 =
\qty{13485.1}{bar}$.

\textbf{14.} $T = (b_2 - b_1)/(a_2 - a_1) = 8193.4/10.412 = \qty{786.92}{K} =
\qty{513.8}{\celsius}$.

\textbf{15.} $P = 11.655\times786.92 - 5291.7 = \qty{3880}{bar} = \qty{3.88}{kbar}$.

\textbf{16.} $(4.01\times786.92 - 3800)/(-0.166) = (3155.5 - 3800)/(-0.166) =
\qty{3882}{bar}$: a mesma pressão, dentro do arredondamento dos dados, como exige a proposição do capítulo.

\textbf{17.} \qty{514}{\celsius} fica a \qty{10}{\celsius} de \qty{504}{\celsius}, dentro
da incerteza de \qty{20}{\celsius}; \qty{3.88}{kbar} fica a \qty{0.13}{kbar} de 3.75,
dentro de 0.25. A concordância é boa.

\textbf{18.} Não: \qty{536}{\celsius} fica a \qty{32}{\celsius} do experimento, fora de sua
incerteza. Desprezar as capacidades térmicas não era o erro principal: as entalpias o são.

\noindent\textbf{Parte IV.}

\textbf{19.} O $\Delta_r H$ de cianita $\to$ andaluzita aumenta \qty{1000}{J}; a $T$
fixo, $P = (T\Delta_r S - \Delta_r H)/\Delta_r V$ cai $1000/0.737 = \qty{1357}{bar}$.

\textbf{20.} $\Delta T = -1357/(a_2 - a_1) = -1357/10.41 = \qty{-130}{K}$: o ponto triplo
se desloca para cerca de \qty{384}{\celsius}.

\textbf{21.} A linha cianita $=$ silimanita cai $1000/0.571 = \qty{1751}{bar}$, e
o ponto triplo se desloca de $+1751/10.41 = \qty{+168}{K}$.

\textbf{22.} Cada entalpia sozinha, variada de $\pm\qty{2.0}{kJ}$, desloca o ponto triplo de
$\pm\qty{260}{K}$ (andaluzita) ou de $\pm\qty{336}{K}$ (silimanita): a calorimetria sozinha admite
pontos triplos espalhados por várias centenas de graus.

\textbf{23.} As entalpias tabeladas não são puramente calorimétricas: dentro de suas
incertezas, elas são ajustadas para que as linhas calculadas reproduzam os experimentos reversos.
Os experimentos fixam o diagrama; as tabelas o transportam, de forma coerente, para outras reações.

\textbf{24.} O ponto triplo calculado a partir das tabelas com propriedades constantes fica a
\textbf{\qty[text-series-to-math=true]{514}{\celsius}}.
\end{solution}
